\section{O algoritmo de travessia}
\label{sec:travessia}

\subsection{Inserção de um lote de arestas}

O ponto de entrada é \code{insertInterfaceElements}, que recebe uma lista de half-edges (uma por aresta; qualquer uma
das duas serve) e as processa na ordem dada (\cref{alg:insert}). Antes de alterar a malha, o método conta os vértices
novos e reserva espaço para eles e para um elemento de quatro half-edges por aresta. Depois, para cada aresta, valida a
entrada, abre o novo elemento a partir do gabarito da \cref{fig:gabarito}, escolhe e aplica o operador em $A$ e em $B$
e, por fim, liga as arestas do meio. A \cref{sec:ibhm-insercao} mostra o efeito de cada passo nos vetores da IBHM.

\begin{algorithm}[ht]
\caption{\code{insertInterfaceElements}$(E)$ \quad (classe \code{IBHMInterfaceElementBuilder})}
\label{alg:insert}
\begin{algorithmic}[1]
\Require lista $E$ de half-edges, uma por aresta a abrir
\State $N\gets$ \code{numberOfHalfEdges}$()$ \Comment{$Of[n_c]$ antes da chamada}
\State $n_v\gets$ \code{computeNumberOfNewVertices}$(E)$
\State \code{reserveSpaceForElements}$(4|E|, |E|)$; \code{reserveSpaceForNodes}$(n_v)$; aumenta
\code{\_isInterfaceElement} e \code{\_inInterfaceElement}
\ForAll{$a\in E$}
  \If{$a\ge N$} \textbf{continue} \Comment{não é half-edge da malha: rejeitada}
  \EndIf
  \State $b\gets\Op[a]$
  \If{$b=\BORDER$ \textbf{ou} $b=\COLLAPSED$} \textbf{continue} \Comment{aresta de borda: rejeitada}
  \EndIf
  \If{$\mathrm{elem}(a)$ \textbf{ou} $\mathrm{elem}(b)$ é elemento de interface} \textbf{continue} \Comment{rejeitada}
  \EndIf
  \State $h_0\gets$ \code{nextAvailableHE}$()$; \quad $h_1,h_2,h_3\gets h_0+1,h_0+2,h_0+3$
  \State $e\gets$ \code{commitElement}$(4)$ \Comment{antes dos operadores: $\nxt/\prv$ dependem de $Of[\,]$}
  \State $\mathit{op}\gets$ \code{retrieveOperator}$(a)$; aplica $\mathit{op}$ com $(a,h_3)$ \Comment{vértice $A$}
  \State $\mathit{op}\gets$ \code{retrieveOperator}$(b)$; aplica $\mathit{op}$ com $(b,h_1)$ \Comment{vértice $B$}
  \State $\Op[h_0]\leftrightarrow b$; \quad $\Op[h_2]\leftrightarrow a$ \Comment{arestas do meio}
  \State \code{\_isInterfaceElement}$[e]\gets$ \textbf{true}
\EndFor
\end{algorithmic}
\end{algorithm}

Alguns detalhes de implementação merecem atenção:
\begin{itemize}
\item \textbf{Rejeições.} Índices que não são half-edges da malha, arestas de borda e arestas que já pertencem a um
elemento de interface (por exemplo, a mesma aresta pedida duas vezes, no mesmo lote ou em chamadas diferentes) são
ignorados com uma mensagem no console. Nada é devolvido a quem chamou (I7), e o espaço reservado para eles fica sem uso.
O teste do índice vem primeiro, porque as consultas seguintes só valem para half-edges de elementos validados. O limite
$N$ é lido no início da chamada: os elementos criados pela própria chamada não contam, e uma half-edge deles na entrada
também é rejeitada. A pré-contagem aplica os mesmos testes (\cref{sec:contagem}), por isso a contagem é exata.
\item \textbf{Validação antecipada do elemento.} \code{commitElement(4)} grava o offset do novo elemento antes de os
operadores rodarem, porque na IBHM $\nxt$ e $\prv$ das half-edges $h_0,\dots,h_3$ dependem dele. Durante os operadores,
as posições $\Op[h_0]$ e $\Op[h_2]$ ainda contêm o \BORDER{} da reserva; nenhuma rotação chega a elas porque nenhuma
half-edge antiga aponta para o novo elemento até o fim da iteração.
\item \textbf{Ordem $A$, depois $B$.} O operador de $B$ é escolhido depois de o operador de $A$ ter sido aplicado. Como
os operadores só mexem na estrela do próprio vértice, isso não altera a decisão em $B$.
\item \textbf{$a$ e $b$ ainda opostas.} Enquanto os operadores trabalham, $a$ e $b$ continuam opostas entre si; é isso
que permite que as rotações em torno de $A$ e de $B$ atravessem a aresta que está sendo aberta (por exemplo, o laço
horário do Open Hole parte de $he$ e passa por $\Op[he]$). Só no final $a$ e $b$ são religadas a $h_2$ e $h_0$.
\item \textbf{Lote ou uma a uma.} Chamar o método uma vez com todas as arestas ou uma vez por aresta produz a mesma
topologia (\cref{sec:verificacao}); a diferença é só a quantidade de reservas de memória.
\end{itemize}

\subsection{O exemplo completo, passo a passo}
\label{sec:exemplo}

A \cref{tab:exemplo} lista o operador escolhido pelo código em cada extremidade, para as 12 arestas do exemplo da
apresentação, na ordem dos slides. Os 24 operadores coincidem exatamente com os anunciados nos slides 15 a 26 da
apresentação~\cite{slides}. A \cref{fig:passos} mostra o estado da malha depois de cada inserção, e a
\cref{fig:final} o resultado final com todos os vértices nomeados.

\begin{table}[ht]
\centering\small
\begin{tabular}{clllc}
\toprule
Aresta & $A$--$B$ (no código) & Operador em $A$ & Operador em $B$ & vértices novos\\
\midrule
1 & $K$--$Q$ & Canônico & Canônico & 1 \\
2 & $K$--$L$ & Split Element & Canônico & 1 \\
3 & $G$--$C$ & Canônico & Canônico & 1 \\
4 & $M$--$H$ & Canônico & Canônico & 1 \\
5 & $S$--$M$ & Canônico & Expand Edge & 2 \\
6 & $L$--$M$ & Expand Edge & Open Hole & 2 \\
7 & $L_{1}$--$G$ & Open Hole & Expand Edge & 2 \\
8 & $K$--$F$ & Open Hole & Canônico & 1 \\
9 & $J$--$K$ & Canônico & Open Hole & 2 \\
10 & $M_{1}$--$N$ & Open Hole & Canônico & 2 \\
11 & $M_{1}$--$I$ & Open Hole & Canônico & 2 \\
12 & $M_{3}$--$T$ & Open Hole & Canônico & 2 \\
\midrule
\multicolumn{4}{r}{total} & 19\\
\bottomrule
\end{tabular}

\caption{Operadores escolhidos pelo código no exemplo (IBHM, malha \code{malha2Hyb.txt}). $X_i$ é a $i$-ésima cópia
criada do vértice $X$. A coluna $A$--$B$ mostra a orientação da half-edge usada (por isso a aresta 1, ``$QK$'' nos
slides, aparece como $K$--$Q$).}
\label{tab:exemplo}
\end{table}

\begin{figure}[p]
\centering
\newcommand{\passo}[2]{\begin{subfigure}[t]{0.32\linewidth}\centering\input{figs/example_step#1.tex}\caption*{\scriptsize #2}\end{subfigure}}
\passo{01}{1: $K$ Canônico, $Q$ Canônico (borda)}\hfill
\passo{02}{2: $K$ Split, $L$ Canônico}\hfill
\passo{03}{3: $G$ Canônico, $C$ Canônico (borda)}\\[3pt]
\passo{04}{4: $M$ Canônico, $H$ Canônico (borda)}\hfill
\passo{05}{5: $S$ Canônico (borda), $M$ Expand}\hfill
\passo{06}{6: $L$ Expand, $M$ Open Hole}\\[3pt]
\passo{07}{7: $L$ Open Hole, $G$ Expand}\hfill
\passo{08}{8: $K$ Open Hole, $F$ Canônico}\hfill
\passo{09}{9: $J$ Canônico (borda), $K$ Open Hole}\\[3pt]
\passo{10}{10: $M$ Open Hole, $N$ Canônico (borda)}\hfill
\passo{11}{11: $M$ Open Hole, $I$ Canônico (borda)}\hfill
\passo{12}{12: $M$ Open Hole, $T$ Canônico (borda)}
\caption{Estado da malha após cada uma das 12 inserções do exemplo (vista explodida). As arestas ainda não inseridas
estão em vermelho e a próxima, em vermelho grosso. Observe a sequência em $M$: colapsado (4), dividido por Expand (5),
buraco triangular (6) que cresce até um hexágono (12). Em $K$ a segunda aresta já chega por Split Element (2), porque
$KL$ e $KQ$ compartilham o elemento $KLRQ$.}
\label{fig:passos}
\end{figure}

\begin{figure}[ht]
\centering
\begin{tikzpicture}[scale=3.200, every node/.style={font=\scriptsize}]
\fill[holefill] (0.880,2.120) -- (1.120,2.120) -- (1.120,1.880) -- (0.880,1.880) -- cycle;
\fill[holefill] (1.880,1.880) -- (2.000,2.170) -- (2.120,1.880) -- cycle;
\fill[holefill] (3.076,1.848) -- (2.880,1.880) -- (2.880,2.120) -- (3.076,2.152) -- (3.152,2.076) -- (3.152,1.924) -- cycle;
\filldraw[cefill] (0.000,0.000) -- (1.000,0.000) -- (1.000,1.000) -- (0.000,1.000) -- cycle;
\filldraw[cefill] (1.000,0.000) -- (1.880,0.120) -- (1.830,1.000) -- (1.000,1.000) -- cycle;
\filldraw[cefill] (2.120,0.120) -- (3.000,0.000) -- (2.830,1.000) -- (2.170,1.000) -- cycle;
\filldraw[cefill] (0.000,1.000) -- (1.000,1.000) -- (0.880,1.880) -- (0.120,1.880) -- cycle;
\filldraw[cefill] (1.000,1.000) -- (1.830,1.000) -- (1.880,1.880) -- (1.120,1.880) -- cycle;
\filldraw[cefill] (2.170,1.000) -- (2.830,1.000) -- (2.880,1.880) -- (2.120,1.880) -- cycle;
\filldraw[cefill] (3.120,1.120) -- (3.848,1.076) -- (3.076,1.848) -- cycle;
\filldraw[cefill] (3.924,1.152) -- (3.880,1.880) -- (3.152,1.924) -- cycle;
\filldraw[cefill] (0.120,2.120) -- (0.880,2.120) -- (0.880,2.880) -- (0.000,3.000) -- cycle;
\filldraw[cefill] (1.120,2.120) -- (2.000,2.170) -- (2.000,3.000) -- (1.120,2.880) -- cycle;
\filldraw[cefill] (2.000,2.170) -- (2.880,2.120) -- (2.880,2.880) -- (2.000,3.000) -- cycle;
\filldraw[cefill] (3.152,2.076) -- (3.880,2.120) -- (3.924,2.848) -- cycle;
\filldraw[cefill] (3.076,2.152) -- (3.848,2.924) -- (3.120,2.880) -- cycle;
\filldraw[iefill] (1.120,2.120) -- (1.120,2.880) -- (0.880,2.880) -- (0.880,2.120) -- cycle;
\filldraw[iefill] (1.120,1.880) -- (1.880,1.880) -- (2.000,2.170) -- (1.120,2.120) -- cycle;
\filldraw[iefill] (1.830,1.000) -- (1.880,0.120) -- (2.120,0.120) -- (2.170,1.000) -- cycle;
\filldraw[iefill] (2.880,1.880) -- (2.830,1.000) -- (3.120,1.120) -- (3.076,1.848) -- cycle;
\filldraw[iefill] (2.880,2.880) -- (2.880,2.120) -- (3.076,2.152) -- (3.120,2.880) -- cycle;
\filldraw[iefill] (2.120,1.880) -- (2.880,1.880) -- (2.880,2.120) -- (2.000,2.170) -- cycle;
\filldraw[iefill] (1.880,1.880) -- (1.830,1.000) -- (2.170,1.000) -- (2.120,1.880) -- cycle;
\filldraw[iefill] (0.880,1.880) -- (0.947,1.046) -- (1.050,1.049) -- (1.120,1.880) -- cycle;
\filldraw[iefill] (0.120,1.880) -- (0.880,1.880) -- (0.880,2.120) -- (0.120,2.120) -- cycle;
\filldraw[iefill] (3.152,1.924) -- (3.880,1.880) -- (3.880,2.120) -- (3.152,2.076) -- cycle;
\filldraw[iefill] (3.076,1.848) -- (3.848,1.076) -- (3.924,1.152) -- (3.152,1.924) -- cycle;
\filldraw[iefill] (3.152,2.076) -- (3.924,2.848) -- (3.848,2.924) -- (3.076,2.152) -- cycle;
\draw[red, line width=1.1pt] (0.947,1.046) -- (1.050,1.049);
\node[vreg, font=\tiny\bfseries, minimum size=0.29cm] at (0.000,0.000) {$A$};
\node[vreg, font=\tiny\bfseries, minimum size=0.29cm] at (1.000,0.000) {$B$};
\node[vie, font=\tiny\bfseries, minimum size=0.29cm] at (2.120,0.120) {$C$};
\node[vreg, font=\tiny\bfseries, minimum size=0.29cm] at (3.000,0.000) {$D$};
\node[vreg, font=\tiny\bfseries, minimum size=0.29cm] at (0.000,1.000) {$E$};
\node[vcol, font=\tiny\bfseries, minimum size=0.29cm] at (1.000,1.000) {$F$};
\node[vie, font=\tiny\bfseries, minimum size=0.29cm] at (1.830,1.000) {$G$};
\node[vie, font=\tiny\bfseries, minimum size=0.29cm] at (3.120,1.120) {$H$};
\node[vie, font=\tiny\bfseries, minimum size=0.29cm] at (3.924,1.152) {$I$};
\node[vie, font=\tiny\bfseries, minimum size=0.29cm] at (0.120,1.880) {$J$};
\node[vie, font=\tiny\bfseries, minimum size=0.29cm] at (0.880,2.120) {$K$};
\node[vie, font=\tiny\bfseries, minimum size=0.29cm] at (2.000,2.170) {$L$};
\node[vie, font=\tiny\bfseries, minimum size=0.29cm] at (2.880,2.120) {$M$};
\node[vie, font=\tiny\bfseries, minimum size=0.29cm] at (3.880,2.120) {$N$};
\node[vreg, font=\tiny\bfseries, minimum size=0.29cm] at (0.000,3.000) {$P$};
\node[vie, font=\tiny\bfseries, minimum size=0.29cm] at (0.880,2.880) {$Q$};
\node[vreg, font=\tiny\bfseries, minimum size=0.29cm] at (2.000,3.000) {$R$};
\node[vie, font=\tiny\bfseries, minimum size=0.29cm] at (2.880,2.880) {$S$};
\node[vie, font=\tiny\bfseries, minimum size=0.29cm] at (3.848,2.924) {$T$};
\node[vie, font=\tiny\bfseries, minimum size=0.29cm] at (1.120,2.880) {$Q_{1}$};
\node[vie, font=\tiny\bfseries, minimum size=0.29cm] at (1.120,2.120) {$K_{1}$};
\node[vie, font=\tiny\bfseries, minimum size=0.29cm] at (1.880,0.120) {$C_{1}$};
\node[vie, font=\tiny\bfseries, minimum size=0.29cm] at (2.830,1.000) {$H_{1}$};
\node[vie, font=\tiny\bfseries, minimum size=0.29cm] at (3.120,2.880) {$S_{1}$};
\node[vie, font=\tiny\bfseries, minimum size=0.29cm] at (3.076,1.848) {$M_{1}$};
\node[vie, font=\tiny\bfseries, minimum size=0.29cm] at (1.880,1.880) {$L_{1}$};
\node[vie, font=\tiny\bfseries, minimum size=0.29cm] at (2.880,1.880) {$M_{2}$};
\node[vie, font=\tiny\bfseries, minimum size=0.29cm] at (2.120,1.880) {$L_{2}$};
\node[vie, font=\tiny\bfseries, minimum size=0.29cm] at (2.170,1.000) {$G_{1}$};
\node[vie, font=\tiny\bfseries, minimum size=0.29cm] at (1.120,1.880) {$K_{2}$};
\node[vie, font=\tiny\bfseries, minimum size=0.29cm] at (0.120,2.120) {$J_{1}$};
\node[vie, font=\tiny\bfseries, minimum size=0.29cm] at (0.880,1.880) {$K_{3}$};
\node[vie, font=\tiny\bfseries, minimum size=0.29cm] at (3.152,2.076) {$M_{3}$};
\node[vie, font=\tiny\bfseries, minimum size=0.29cm] at (3.880,1.880) {$N_{1}$};
\node[vie, font=\tiny\bfseries, minimum size=0.29cm] at (3.152,1.924) {$M_{4}$};
\node[vie, font=\tiny\bfseries, minimum size=0.29cm] at (3.848,1.076) {$I_{1}$};
\node[vie, font=\tiny\bfseries, minimum size=0.29cm] at (3.076,2.152) {$M_{5}$};
\node[vie, font=\tiny\bfseries, minimum size=0.29cm] at (3.924,2.848) {$T_{1}$};

\end{tikzpicture}

\caption{Resultado final do exemplo, com os nomes de todos os vértices. Os vértices $K$, $L$ e $M$ terminam com 4, 3 e 6
cópias e com buracos de 4, 3 e 6 lados; $G$ (duas arestas) termina com duas cópias coladas; $F$ (uma aresta, interno)
continua com uma aresta colapsada (contorno vermelho); os vértices de borda $C$, $H$, $I$, $J$, $N$, $Q$, $S$ e $T$
foram duplicados uma vez. A posição das cópias é apenas ilustrativa: no código cada cópia tem exatamente as
coordenadas do seu vértice de origem.}
\label{fig:final}
\end{figure}

\subsection{Propriedades observadas}

\paragraph{Independência da ordem.} O resultado final não depende da ordem de inserção, da half-edge escolhida para
cada aresta, nem de as arestas virem em lote ou uma a uma. A razão é que a topologia final é determinada pelo
\emph{conjunto} de arestas: as arestas selecionadas cortam a estrela de cada vértice em setores, e cada setor termina com
uma cópia própria (salvo no vértice interno com uma única aresta, que fica colapsado). Os operadores sempre dividem a
estrela exatamente ao longo das arestas selecionadas, de modo que qualquer ordem leva à mesma partição. Verificamos isso
para três ordens do exemplo (a dos slides, a de \code{main.cpp} e o lote inteiro em uma única chamada), para outras
21 ordens nos testes automatizados (\cref{sec:testes}) e, nos 659 casos aleatórios que terminaram
(\cref{sec:verificacao}), para três ordens aleatórias e uma inserção em lote.

\paragraph{Custo.} Cada aresta custa $O\big((\deg A+\deg B)\log n_c\big)$, dominado pelas rotações (\code{isBorder} no
Canônico, reindexação no Expand e no Open Hole); o fator $\log n_c$ vem da busca binária em $Of[\,]$ que toda consulta
$\nxt$/$\prv$ faz. A pré-contagem também chama \code{isBorder}, uma vez por vértice distinto. O custo total é proporcional ao
número de arestas inseridas para malhas de valência limitada, e só depende do tamanho da malha pelo fator logarítmico.

\paragraph{Memória.} O número de vértices novos é exatamente o previsto, inclusive com arestas rejeitadas ou repetidas
no lote. É reservado um elemento por aresta do lote, e o espaço das arestas rejeitadas fica sem uso.

\paragraph{O exemplo no \code{main.cpp}.} Hoje o \code{main.cpp} executa este construtor sobre \code{malha2Hyb.txt},
com as mesmas 12 arestas dos slides em outra ordem (9, 1, 2, 8, 6, 7, 3, 10, 4, 5, 11, 12, na numeração dos slides); o
resultado é o mesmo da \cref{fig:final}.
